\chapter{The DQC Dialect: Designing a Domain-Specific IR in MLIR}
\label{ch:mlir_compiler_architecture}

\begin{quote}
\textit{``Monolithic quantum compilers treat circuits as flat arrays of gates. To compile distributed quantum computers, we must elevate quantum semantics into a first-class Multi-Level Intermediate Representation with formal types, verifiable invariants, and progressive lowering.''}
\end{quote}

% ======================================================================
% OPENING NARRATIVE
% ======================================================================
\section*{Chapter Overview: The Compiler Abstraction Crisis in Quantum Computing}

For the past decade, quantum compiler engineering has suffered from what we call the \textbf{monolithic abstraction crisis}. Early quantum software frameworks (such as Qiskit, Cirq, ProjectQ, and Quil) represented quantum programs as either:
\begin{enumerate}
    \item Flat sequential lists of Python objects (\texttt{circuit.append(CNOT(0, 1))}), or
    \item Global monolithic Directed Acyclic Graphs (DAGs) where nodes represent gates and directed edges represent qubit dependencies.
\end{enumerate}

While this flat representation was adequate for 5-qubit single-chip prototypes, it completely collapses when targeting distributed multi-QPU systems. Consider the simultaneous compiler responsibilities in a distributed quantum cluster:
\begin{itemize}
    \item \textbf{Quantum Unitary Layer:} Canceling redundant gate sequences ($H \cdot H \to \mathbb{I}$, rotation merging, KAK decomposition);
    \item \textbf{Spatial Graph Layer:} Partitioning qubits across cryostats and optimizing hypergraph cut sizes;
    \item \textbf{Classical-Quantum Hybrid Control Layer:} Modeling mid-circuit measurements, branch conditions, and microsecond real-time feedforward fixup logic;
    \item \textbf{Asynchronous Network Layer:} Managing distributed EPR generation queues, photon detector heralds, and MPI/RDMA inter-rack network packets;
    \item \textbf{Physical Timing Layer:} Scheduling pulse waveforms against $T_1$ relaxation and $T_2^*$ dephasing deadlines in nanosecond resolution.
\end{itemize}

Attempting to express all five levels of abstraction within a single flat Python DAG results in brittle, unmaintainable compilers. What is needed is a \textbf{Multi-Level Intermediate Representation (MLIR)}---a compiler infrastructure where multiple distinct dialects coexist, each tailored to a specific level of abstraction, progressively lowering programs from high-level algorithms to distributed network primitives.

In this chapter, we design and implement the \keyterm{DQC Dialect} within LLVM's MLIR framework. We detail TableGen specifications, formalize the linear type system reconciling Static Single Assignment (SSA) with the No-Cloning Theorem, define declarative rewrite patterns in C++, and construct verifiably sound progressive lowering pipelines.

% ==============================================================================
\section{Why MLIR for Distributed Quantum Compilation?}
\label{sec:why_mlir}
% ==============================================================================

MLIR (Multi-Level Intermediate Representation) is an open-source compiler infrastructure originated within the LLVM project. Unlike traditional compilers that force all semantics into a single rigid intermediate representation (like LLVM IR), MLIR allows compiler architects to define modular, domain-specific vocabularies called \keyterm{Dialects}.

\begin{figure}[htbp]
    \centering
    \begin{tikzpicture}[scale=0.92, >=stealth]
        % Level 1: High Level AST
        \draw[rounded corners=2mm, fill=blue!10, draw=darkblue, line width=1pt] (-4.5, 4.2) rectangle (4.5, 5.1);
        \node[font=\bfseries\color{darkblue}] at (0, 4.65) {High-Level Quantum Algorithms (Qiskit, OpenQASM 3.0, Cirq)};
        
        % Level 2: DQC Monolithic Dialect
        \draw[rounded corners=2mm, fill=blue!20, draw=darkblue, line width=1pt] (-4.5, 2.7) rectangle (4.5, 3.6);
        \node[font=\bfseries\color{darkblue}] at (0, 3.15) {\texttt{dqc} Dialect (Monolithic IR: Gates, Unitary Traits)};
        
        % Level 3: Partitioned DQC Dialect
        \draw[rounded corners=2mm, fill=compilerteal!20, draw=compilerteal, line width=1pt] (-4.5, 1.2) rectangle (4.5, 2.1);
        \node[font=\bfseries\color{compilerteal}] at (0, 1.65) {\texttt{dqc.dist} Dialect (QPU Placement, Non-Local Telegates)};
        
        % Level 4: Split Lowering (Quantum + Classical Communication)
        \draw[rounded corners=2mm, fill=quantumviolet!15, draw=quantumviolet, line width=1pt] (-5.8, -0.7) rectangle (-0.5, 0.3);
        \node[font=\small\bfseries\color{quantumviolet}, align=center] at (-3.15, -0.2) {QIR / OpenQASM Dialect\\(Local Pulse Schedules)};
        
        \draw[rounded corners=2mm, fill=quantumviolet!15, draw=quantumviolet, line width=1pt] (0.5, -0.7) rectangle (5.8, 0.3);
        \node[font=\small\bfseries\color{quantumviolet}, align=center] at (3.15, -0.2) {Async / MPI / LLVM Dialect\\(Classical Sync \& Teleportation)};
        
        % Level 5: Target Hardware
        \draw[rounded corners=2mm, fill=gray!20, draw=black, line width=1.2pt] (-4.5, -2.4) rectangle (4.5, -1.4);
        \node[font=\bfseries, align=center] at (0, -1.9) {Distributed Quantum Hardware Controller / FPGA Cluster};
        
        % Arrows
        \draw[->, line width=1.5pt, color=darkblue] (0, 4.2) -- node[right, font=\tiny\bfseries] {Ingestion} (0, 3.6);
        \draw[->, line width=1.5pt, color=darkblue] (0, 2.7) -- node[right, font=\tiny\bfseries] {Pass 2: Hypergraph Partitioning} (0, 2.1);
        \draw[->, line width=1.5pt, color=compilerteal] (-2, 1.2) -- (-3.15, 0.3);
        \draw[->, line width=1.5pt, color=compilerteal] (2, 1.2) -- (3.15, 0.3);
        \draw[->, line width=1.5pt, color=quantumviolet] (-3.15, -0.7) -- (-2, -1.4);
        \draw[->, line width=1.5pt, color=quantumviolet] (3.15, -0.7) -- (2, -1.4);
    \end{tikzpicture}
    \caption{The Multi-Level IR Progressive Lowering Pipeline in DQC. The compiler preserves high-level algorithmic intent, progressively transforming monolithic quantum operations into partitioned distributed primitives, then splitting into parallel hardware execution streams (QIR for local qubits, MPI/LLVM for inter-refrigerator network messaging).}
    \label{fig:mlir_lowering_hierarchy_ch07}
\end{figure}

The core advantages of MLIR for distributed quantum compilation include:
\begin{enumerate}
    \item \textbf{Progressive Lowering:} Compilers transform code through distinct, verifiable semantic stages rather than attempting an all-at-once translation from Python to FPGA microcode.
    \item \textbf{Declarative Op Definitions (ODS):} Operations, types, and invariants are written in TableGen (\texttt{.td}), from which C++ classes, verifiers, and serialization parsers are automatically synthesized.
    \item \textbf{Declarative Rewrite Rules (DRR):} Quantum algebraic equivalence patterns ($H \cdot H \to \mathbb{I}$, $X \cdot Z \to -iY$) are expressed in high-level pattern matching syntax.
    \item \textbf{Unified Classical-Quantum Coexistence:} Classical integer loop indices, memory pointers, and quantum state handles live seamlessly in the same Basic Block and Control Flow Graph (CFG).
\end{enumerate}

% ==============================================================================
\section{Formal Linear Type System and No-Cloning Enforcement}
\label{sec:linear_type_system}
% ==============================================================================

In classical MLIR dialects (e.g., \texttt{arith} or \texttt{math}), SSA values represent read-only immutable data that can be used arbitrarily many times:
\begin{lstlisting}[language=MLIR]
%c = arith.addi %a, %b : i32
%d = arith.muli %c, %c : i32 // Valid: %c is used twice!
\end{lstlisting}

In quantum computing, this is physically illegal. If \texttt{\%q} represents a physical quantum state $\ket{\psi}$, using \texttt{\%q} as an operand in two different gates simultaneously would clone the state, violating the No-Cloning Theorem!

To resolve this, the DQC dialect enforces a \textbf{Linear Type System}: every quantum operation consumes its input SSA qubit values and produces new, updated SSA qubit handles.

\subsection{Typing Judgment Rules and Operational Semantics}
Let $\Gamma$ be the typing context mapping SSA identifiers to types. We write $\Gamma_1 \circ \Gamma_2 = \Gamma$ for linear context splitting where the variables in $\Gamma$ are partitioned disjointly into $\Gamma_1$ and $\Gamma_2$.

\begin{definitionbox}{Linear Quantum Typing Rules}
\begin{align}
    \text{(Variable)} \quad & \frac{}{\{x : \tau\} \vdash x : \tau} \label{eq:type_var_ch07} \\[1ex]
    \text{(1Q Gate)} \quad & \frac{\Gamma \vdash q : \texttt{!dqc.qubit}}{\Gamma \vdash \texttt{dqc.rot}(q, \theta, \phi, \lambda) : \texttt{!dqc.qubit}} \label{eq:type_1q_ch07} \\[1ex]
    \text{(2Q Gate)} \quad & \frac{\Gamma_1 \vdash q_1 : \texttt{!dqc.qubit} \quad \Gamma_2 \vdash q_2 : \texttt{!dqc.qubit}}{\Gamma_1 \circ \Gamma_2 \vdash \texttt{dqc.cnot}(q_1, q_2) : \texttt{!dqc.qubit} \times \texttt{!dqc.qubit}} \label{eq:type_2q_ch07} \\[1ex]
    \text{(Telegate)} \quad & \frac{\Gamma_1 \vdash q_1 : \texttt{!dqc.qubit} \quad \Gamma_2 \vdash q_2 : \texttt{!dqc.qubit} \quad \Gamma_3 \vdash e : \texttt{!dqc.epr\_pair}}{\Gamma_1 \circ \Gamma_2 \circ \Gamma_3 \vdash \texttt{dqc.telegate}(q_1, q_2, e) : \texttt{!dqc.qubit} \times \texttt{!dqc.qubit}} \label{eq:type_telegate_ch07}
\end{align}
\end{definitionbox}

\begin{theorembox}{Type Safety and No-Cloning Soundness}
If an MLIR module $\mathcal{M}$ is well-typed under the linear DQC typing rules ($\vdash \mathcal{M} : \text{valid}$), then no physical quantum state handle is cloned, duplicated, or consumed concurrently across multiple execution paths.
\end{theorembox}

\begin{proof}
By induction on the derivation tree. Axiom~\eqref{eq:type_var_ch07} guarantees that each variable in $\Gamma$ can be referenced at most once. In rules~\eqref{eq:type_2q_ch07} and \eqref{eq:type_telegate_ch07}, the linear context splitting $\Gamma = \Gamma_1 \circ \Gamma_2 \circ \Gamma_3$ forces $\text{dom}(\Gamma_i) \cap \text{dom}(\Gamma_j) = \emptyset$, preventing the same SSA qubit identifier from appearing simultaneously in distinct operands.
\end{proof}

% ==============================================================================
\section{Dialect Definition and TableGen Specifications}
\label{sec:tablegen_spec}
% ==============================================================================

In MLIR, dialects, types, and operations are defined declaratively in TableGen (\texttt{.td}) files.

\subsection{Root Dialect Declaration}
\begin{lstlisting}[language=MLIR, caption={TableGen root definition of the DQC Dialect (DQCDialect.td)}]
// DQCDialect.td - Root Dialect Definition
def DQC_Dialect : Dialect {
  let name = "dqc";
  let summary = "Distributed Quantum Computing Dialect";
  let description = [{
    The DQC dialect provides first-class primitives for modeling, 
    partitioning, synthesizing, and scheduling distributed quantum circuits 
    across heterogeneous multi-QPU networks.
  }];
  let cppNamespace = "::mlir::dqc";
  let useDefaultTypePrinterParser = 1;
}
\end{lstlisting}

\subsection{Type System Definitions}
\begin{lstlisting}[language=MLIR, caption={TableGen Type Definitions for DQC (DQCTypes.td)}]
// Custom Types in DQC Dialect
def DQC_QubitType : TypeDef<DQC_Dialect, "Qubit"> {
  let mnemonic = "qubit";
  let summary = "Linear SSA handle to a physical or logical quantum bit";
}

def DQC_EPRPairType : TypeDef<DQC_Dialect, "EPRPair"> {
  let mnemonic = "epr_pair";
  let summary = "Handle to a pre-shared entangled Bell pair |Phi+>";
}

def DQC_QNodeType : TypeDef<DQC_Dialect, "QNode"> {
  let mnemonic = "qnode";
  let summary = "Physical QPU node identifier in distributed cluster";
}
\end{lstlisting}

\subsection{TableGen Operation Definitions}
\begin{lstlisting}[language=MLIR, caption={Core TableGen Operations in DQC (DQCOps.td)}]
// Base class for DQC operations
class DQC_Op<string mnemonic, list<Trait> traits = []> :
    Op<DQC_Dialect, mnemonic, traits>;

// 1. Qubit Allocation on Specific QPU
def DQC_AllocQubitOp : DQC_Op<"alloc_qubit"> {
  let summary = "Allocate a new physical qubit in |0> state on a specific QPU";
  let arguments = (ins OptionalAttr<I32Attr>:$qpu_id);
  let results = (outs DQC_QubitType:$qubit);
  let assemblyFormat = "(`on` $qpu_id^)? attr-dict `:` type($qubit)";
}

// 2. Single-Qubit Rotation Gate
def DQC_RotOp : DQC_Op<"rot"> {
  let summary = "Arbitrary single-qubit rotation R(theta, phi, lambda)";
  let arguments = (ins DQC_QubitType:$in_qubit, F64Attr:$theta, F64Attr:$phi, F64Attr:$lambda);
  let results = (outs DQC_QubitType:$out_qubit);
  let assemblyFormat = "$in_qubit `[` $theta `,` $phi `,` $lambda `]` attr-dict `:` type($in_qubit)";
}

// 3. Single-Qubit Hadamard Gate
def DQC_HOp : DQC_Op<"h"> {
  let summary = "Single-qubit Hadamard superposition gate";
  let arguments = (ins DQC_QubitType:$in_qubit);
  let results = (outs DQC_QubitType:$out_qubit);
  let assemblyFormat = "$in_qubit attr-dict `:` type($in_qubit)";
}

// 4. Two-Qubit CNOT Gate
def DQC_CNOTOp : DQC_Op<"cnot"> {
  let summary = "Controlled-NOT gate between control and target qubits";
  let arguments = (ins DQC_QubitType:$control, DQC_QubitType:$target);
  let results = (outs DQC_QubitType:$out_control, DQC_QubitType:$out_target);
  let assemblyFormat = "$control `,` $target attr-dict `:` type($control) `,` type($target)";
  let hasVerifier = 1;
}

// 5. Pre-Shared EPR Pair Preparation
def DQC_PrepareEprOp : DQC_Op<"prepare_epr"> {
  let summary = "Generate an entangled Bell pair between two QPU nodes";
  let arguments = (ins I32Attr:$src_qpu, I32Attr:$dst_qpu, F64Attr:$target_fidelity);
  let results = (outs DQC_EPRPairType:$epr);
  let assemblyFormat = "$src_qpu `<->` $dst_qpu `fid` $target_fidelity attr-dict `:` type($epr)";
  let hasVerifier = 1;
}

// 6. Non-Local Telegate Primitive
def DQC_TelegateOp : DQC_Op<"telegate"> {
  let summary = "Synthesized non-local quantum gate using an EPR pair";
  let arguments = (ins DQC_QubitType:$qubit_src, 
                       DQC_QubitType:$qubit_dst, 
                       DQC_EPRPairType:$epr,
                       StrAttr:$gate_kind);
  let results = (outs DQC_QubitType:$out_src, DQC_QubitType:$out_dst);
  let hasVerifier = 1;
}

// 7. Mid-Circuit Measurement
def DQC_MeasureOp : DQC_Op<"measure"> {
  let summary = "Dispersive measurement of a physical qubit in the computational basis";
  let arguments = (ins DQC_QubitType:$qubit);
  let results = (outs I1:$result);
  let assemblyFormat = "$qubit attr-dict `:` type($qubit)";
}
\end{lstlisting}

% ==============================================================================
\section{C++ Custom Verifiers and Dialect Invariants}
\label{sec:custom_verifiers}
% ==============================================================================

MLIR allows attaching custom C++ verifiers to guarantee semantic invariants at compile time.

\begin{lstlisting}[language=C++, caption={C++ Verifiers for DQC Operations (DQCOps.cpp)}]
#include "dqc/DQCDialect.h"
#include "dqc/DQCOps.h"
#include "mlir/IR/OpImplementation.h"

namespace mlir::dqc {

::mlir::LogicalResult CNOTOp::verify() {
  if (getControl() == getTarget()) {
    return emitOpError("Control and Target qubits must be distinct SSA values (No-Cloning violation).");
  }
  return ::mlir::success();
}

::mlir::LogicalResult TelegateOp::verify() {
  auto srcQubit = getQubitSrc();
  auto dstQubit = getQubitDst();

  // Invariant 1: Source and Target qubits must be distinct
  if (srcQubit == dstQubit) {
    return emitOpError("Source and target qubits in a telegate must be distinct SSA values.");
  }

  // Invariant 2: Gate kind attribute must be a valid supported unitary
  llvm::StringRef kind = getGateKind();
  if (kind != "CNOT" && kind != "CZ" && kind != "SWAP" && kind != "CP") {
    return emitOpError("Unsupported non-local gate kind: '") << kind << "'.";
  }

  return ::mlir::success();
}

::mlir::LogicalResult PrepareEprOp::verify() {
  if (getSrcQpu() == getDstQpu()) {
    return emitOpError("EPR pair generation requires two distinct QPU IDs.");
  }
  if (getTargetFidelity() <= 0.5 || getTargetFidelity() > 1.0) {
    return emitOpError("Target EPR fidelity must lie in the entangled range (0.5, 1.0].");
  }
  return ::mlir::success();
}

} // namespace mlir::dqc
\end{lstlisting}

\subsection{Declarative Rewrite Rules (DRR) for Quantum Algebraic Simplifications}
\label{sec:drr_quantum_patterns}

In MLIR, quantum algebraic simplifications are written declaratively in TableGen pattern rules (\texttt{DQCPatterns.td}). During the canonicalization pass, the pattern-match-and-rewrite engine automatically identifies and collapses algebraic redundancies, as detailed in Listing~\ref{lst:drr_patterns_td}.

\begin{lstlisting}[language=MLIR, caption={Declarative Quantum Rewrite Patterns (DQCPatterns.td)}, label={lst:drr_patterns_td}]
// DQCPatterns.td - Declarative Algebraic Quantum Rewrite Rules
include "dqc/DQCOps.td"
include "mlir/IR/PatternBase.td"

// 1. Self-Inverse Involutions: H * H = Identity
def EliminateAdjacentHadamards : Pat<
  (DQC_HOp (DQC_HOp $qubit)),
  (replaceWithValue $qubit)
>;

// 2. Self-Inverse Involutions: X * X = Identity, Z * Z = Identity
def EliminateAdjacentPauliX : Pat<
  (DQC_XOp (DQC_XOp $qubit)),
  (replaceWithValue $qubit)
>;

def EliminateAdjacentPauliZ : Pat<
  (DQC_ZOp (DQC_ZOp $qubit)),
  (replaceWithValue $qubit)
>;

// 3. Continuous Rotation Merging: Rz(theta1) * Rz(theta2) -> Rz(theta1 + theta2)
def MergeAdjacentRzRotations : Pat<
  (DQC_RzOp (DQC_RzOp $qubit, $theta1), $theta2),
  (DQC_RzOp $qubit, (Arith_AddFOp $theta1, $theta2))
>;

// 4. Telegate Commutation: Remote CZ commutes with Local Z
def CommuteRemoteCZWithLocalZ : Pat<
  (DQC_TelegateOp (DQC_ZOp $qA), $qB, $epr, ConstantStrAttr<"CZ">),
  (DQC_ZOp (DQC_TelegateOp $qA, $qB, $epr, ConstantStrAttr<"CZ">))
>;
\end{lstlisting}

\subsection{Complete C++ Implementation of the DQC Linear Verifier Pass}
\label{sec:cpp_linear_verifier_pass}

To guarantee that no quantum circuit violating the No-Cloning theorem or executing illegal use-after-consume SSA operations passes to downstream code generation, the DQC compiler runs a dedicated verification pass (\texttt{DQCLinearVerifierPass.cpp}), detailed in Listing~\ref{lst:linear_verifier_cpp}.

\begin{lstlisting}[language=C++, caption={Complete C++ Linear SSA Verifier Pass (DQCLinearVerifierPass.cpp)}, label={lst:linear_verifier_cpp}]
#include "dqc/DQCDialect.h"
#include "dqc/DQCOps.h"
#include "dqc/DQCPasses.h"
#include "mlir/Pass/Pass.h"
#include "llvm/ADT/DenseSet.h"

namespace mlir::dqc {

#define GEN_PASS_DEF_DQCLINEARVERIFIERPASS
#include "dqc/DQCPasses.h.inc"

struct DQCLinearVerifierPass 
    : public impl::DQCLinearVerifierPassBase<DQCLinearVerifierPass> {

  void runOnOperation() override {
    ModuleOp module = getOperation();
    bool verificationFailed = false;

    module.walk([&](func::FuncOp funcOp) {
      for (Block &block : funcOp.getBody()) {
        llvm::DenseSet<Value> consumedValues;

        for (Operation &op : block) {
          // Check all operands for linear consumption violations
          for (Value operand : op.getOperands()) {
            if (llvm::isa<QubitType, EPRPairType>(operand.getType())) {
              if (consumedValues.contains(operand)) {
                op.emitOpError("Linear type violation: Qubit/EPR SSA value used after consumption.")
                    .attachNote(operand.getLoc(), "Operand first consumed here");
                verificationFailed = true;
                return WalkResult::interrupt();
              }
              consumedValues.insert(operand);
            }
          }
        }
      }
      return WalkResult::advance();
    });

    if (verificationFailed) {
      signalPassFailure();
    }
  }
};

std::unique_ptr<Pass> createDQCLinearVerifierPass() {
  return std::make_unique<DQCLinearVerifierPass>();
}

} // namespace mlir::dqc
\end{lstlisting}

\begin{algorithmbox}{Linear Type SSA Verification Algorithm}
\textbf{Input:} MLIR Block $\mathcal{B}$ containing quantum operations.\\
\textbf{Output:} \texttt{LogicalResult::success()} or \texttt{emitError()}.

\begin{enumerate}[leftmargin=1.5em]
    \item Initialize consumed value set $\mathcal{S}_{\text{consumed}} \leftarrow \emptyset$.
    \item \textbf{For each} operation $op \in \mathcal{B}$ \textbf{do}:
    \begin{itemize}
        \item \textbf{For each} operand $v \in op.\text{getOperands()}$:
        \begin{itemize}
            \item If $v.\text{getType()}$ is \texttt{DQC\_QubitType} or \texttt{DQC\_EPRPairType}:
            \begin{itemize}
                \item If $v \in \mathcal{S}_{\text{consumed}}$: \textbf{return} \texttt{emitOpError("Use-after-consume violation on qubit SSA value")}.
                \item Insert $v$ into $\mathcal{S}_{\text{consumed}}$.
            \end{itemize}
        \end{itemize}
    \end{itemize}
    \item \textbf{Return} \texttt{success()}.
\end{enumerate}
\end{algorithmbox}

\subsection{Complete Program Transformation Listing}
To see the DQC dialect in action, consider a 2-qubit Bell state circuit compiled from monolithic unpartitioned IR to distributed QPU-annotated IR:

\begin{lstlisting}[language=MLIR, caption={Monolithic DQC IR before partitioning}]
module {
  func.func @bell_state_monolithic() -> (!dqc.qubit, !dqc.qubit) {
    %q0 = dqc.alloc_qubit : !dqc.qubit
    %q1 = dqc.alloc_qubit : !dqc.qubit
    %q0_h = dqc.h(%q0) : !dqc.qubit
    %q0_out, %q1_out = dqc.cnot %q0_h, %q1 : !dqc.qubit, !dqc.qubit
    return %q0_out, %q1_out : !dqc.qubit, !dqc.qubit
  }
}
\end{lstlisting}

\begin{lstlisting}[language=MLIR, caption={Partitioned Distributed DQC IR after Pass 2 \& 3}]
module {
  func.func @bell_state_distributed() -> (!dqc.qubit, !dqc.qubit) {
    // Qubit 0 on Cryostat 0, Qubit 1 on Cryostat 1
    %q0 = dqc.alloc_qubit on 0 : !dqc.qubit
    %q1 = dqc.alloc_qubit on 1 : !dqc.qubit
    %q0_h = dqc.h(%q0) : !dqc.qubit
    
    // Allocate pre-shared EPR Bell pair between QPU 0 and QPU 1
    %epr = dqc.prepare_epr 0 <-> 1 fid 0.98 : !dqc.epr_pair
    
    // Synthesized Non-Local Tele-CNOT Gate
    %q0_out, %q1_out = dqc.telegate %q0_h, %q1, %epr {gate_kind = "CNOT"} : !dqc.qubit, !dqc.qubit
    return %q0_out, %q1_out : !dqc.qubit, !dqc.qubit
  }
}
\end{lstlisting}

% ==============================================================================
% CHAPTER SUMMARY
% ==============================================================================

\begin{summarybox}{Key Invariants of Chapter 7}
\begin{enumerate}
    \item \textbf{Progressive Lowering:} MLIR enables modular multi-level lowering from algorithmic DAGs to partitioned distributed pulses and network sync messages.
    \item \textbf{Linear SSA Types:} The \texttt{!dqc.qubit} type enforces linear value consumption, guaranteeing compile-time adherence to the No-Cloning theorem.
    \item \textbf{TableGen Framework:} Declarative ODS definitions synthesize robust C++ AST classes, serialization parsers, and custom verification hooks.
    \item \textbf{EPR First-Class Citizens:} Entanglement pairs (\texttt{!dqc.epr\_pair}) are explicit first-class compiler types subject to strict provenance and purification constraints.
\end{enumerate}
\end{summarybox}

% ==============================================================================
% EXERCISES
% ==============================================================================

\section*{Exercises}
\addcontentsline{toc}{section}{Exercises}

\begin{enumerate}[label=\textbf{7.\arabic*.}]
    \item \textbf{Linear Type Invariants.} Explain why classical SSA allows multiple uses of a value \texttt{\%x}, whereas quantum SSA requires single-use consumption.
    \item \textbf{TableGen Op Definition.} Write the complete TableGen definition for a 3-qubit Toffoli gate \texttt{dqc.toffoli} taking 2 controls and 1 target.
    \item \textbf{Verifier Implementation.} Write a C++ verifier for a \texttt{dqc.purify} operation ensuring that input EPR pairs have matching node IDs.
    \item \textbf{Declarative Rewrite Rule (DRR).} Write a TableGen pattern matching rule that replaces two adjacent Hadamard gates ($H \cdot H$) on the same qubit with a no-op identity wire.
    \item \textbf{Dialect Lowering Design.} Design the lowering pass converting \texttt{dqc.prepare\_epr} into asynchronous MPI message calls (\texttt{MPI\_Isend}, \texttt{MPI\_Irecv}) and optical pulse triggers.
    \item \textbf{Typing Rule Derivation.} Using the linear typing rules, construct the formal typing derivation tree for a 3-qubit GHZ state preparation circuit in the DQC dialect.
    \item \textbf{Memory Effect Trait.} Explain how the \texttt{MemoryEffectOpInterface} trait in MLIR is parameterized to model decoherence side-effects on idle qubits during long communication waits.
    \item \textbf{Custom Dialect Attribute.} Design a custom TableGen attribute \texttt{DQC\_NoiseModelAttr} that stores $T_1, T_2^*$, and readout error probabilities for a given QPU.
\end{enumerate}
